\documentclass[12pt]{article}
\usepackage[utf8]{inputenc}
\usepackage[spanish]{babel}
\usepackage[margin=2.5cm]{geometry}
\usepackage{booktabs}
\usepackage{array}

\begin{document}

\begin{center}
    \textbf{UNIVERSIDAD TÉCNICA ESTATAL DE QUEVEDO}\\[4pt]
    \textbf{INFORME DE LA EXPOSICIÓN 2DO CORTE: DEFENSA DEL PFC}\\[4pt]
    \textbf{Nombre: Jeremy Ruperto Gaibor Rodríguez}
\end{center}

\vspace{6pt}

Durante la actividad se observaron las exposiciones de tres grupos. Para realizar el informe se tomó en cuenta el porcentaje de participación de cada integrante, la claridad de la explicación, la relación entre los temas presentados, la ampliación del contenido de las diapositivas y las evidencias mostradas sobre la implementación de cada proyecto.

Los porcentajes de cumplimiento son aproximados y se establecieron de acuerdo con lo observado durante las exposiciones.

\section*{Grupo Pacheco, Carpio, Cando y Álvarez}

Este grupo obtuvo un cumplimiento aproximado del \textbf{95\%} y, de acuerdo con lo observado, fue el que realizó la exposición más completa.

Durante la presentación explicaron el problema del proyecto, la arquitectura utilizada, las decisiones técnicas y los avances realizados. También mostraron evidencias relacionadas con la fragmentación por fecha, el clúster de CockroachDB, las métricas de Prometheus, las pruebas con Testcontainers y el procesamiento de datos mediante Spark. La exposición terminó con una demostración del funcionamiento del sistema y con los puntos que todavía quedaban pendientes.

La información se presentó siguiendo una secuencia entendible, desde la explicación del problema hasta los resultados. Además, no se limitaron al contenido de las diapositivas, ya que mostraron pruebas, métricas y una demostración práctica.

En este grupo no se anotó qué tema explicó individualmente cada integrante, por lo que solo se registra el porcentaje de participación.

\begin{center}
\begin{tabular}{lc}
\toprule
\textbf{Integrante} & \textbf{Participación} \\
\midrule
Pacheco & 47\% \\
Carpio & 30\% \\
Cando & 14\% \\
Álvarez & 9\% \\
\midrule
\textbf{Total} & \textbf{100\%} \\
\bottomrule
\end{tabular}
\end{center}

Aunque la distribución del tiempo no fue equilibrada, el grupo presentó suficiente evidencia del trabajo realizado y explicó tanto la parte técnica como el funcionamiento del proyecto.



\end{document}